% ch-policy.tex -- what a version number promises, and what may change between versions
\chapter{Versions and updates}
\label{ch:policy}

Yuzen is at version 1.0. This chapter says what the version number means, what
you may rely on between one release and the next, and --- the part that matters
most in a solver --- what happens when a change moves a number. It is a
statement of policy rather than a description of machinery, so where the code
enforces something the text says so, and where it does not, it says that too.

\textbf{Version 1.0 is the first public release} (\cref{ch:revisions}).

\section{The version number}
\label{sec:policy:scheme}

The version is \code{MAJOR.MINOR.PATCH}, defined by three macros in
\file{yuzen/src/Version.h} and formatted by \code{C\_Yuzen::Version()}. Each
level is a promise to you.

\begin{description}

  \item[PATCH] Bug fixes and documentation. \textbf{Existing input files keep
    working, unchanged.} Numerical results may change \emph{only} where the
    previous behavior was wrong, and any such change is listed in
    \cref{ch:revisions} together with the size of the difference and the case
    it was measured on. A patch release that moves a number without saying so
    is a defect in the release, not a permissible one.

  \item[MINOR] New features: new input keys, new allowed values for an existing
    key, new physics terms, new output columns. \textbf{Existing input files
    keep working} --- a key that gains a new allowed value, or an optional key
    that did not exist before, breaks nothing, because an absent key takes its
    default. Results of existing cases do not change, for the reason set out in
    \cref{sec:repro:order}: a new term is appended to the registration list and
    is disabled by default, so a case that does not use it is bit-identical.

  \item[MAJOR] A change that can invalidate existing input files or existing
    results: a key whose meaning changes, a default that changes, a
    discretization that is replaced. A major release is where such a change is
    allowed to happen, and where you should expect to re-run anything you care
    about.

\end{description}

Then there is the rule the code actually implements.
\code{C\_Yuzen::VersionSeries()} returns \code{MAJOR.MINOR} only, and
\code{Input::ReadRoot()} in \file{yuzen/src/IO/Input.cpp} checks
\key{SolverName} against a list of accepted series. So \key{SolverName} in an input file names
\code{MAJOR.MINOR} and never the patch level: at version 1.0.0 the list holds
the single entry \val{"Yuzen-1.0"}, and anything else ends the run with
\code{Invalid value. SolverName must be one of [...]}, the list itself printed in
the message.

\textbf{A patch release therefore never requires editing an input file}, which
is the point of the rule.

\begin{note}[How a minor release keeps older input files working]
  The check is against a \emph{list} of accepted series, built in
  \code{Input::ReadRoot()}, which at 1.0 holds the single entry
  \val{"Yuzen-1.0"}. A minor release that does not change how any key is
  interpreted \textbf{adds} its new series to that list and leaves the old
  entries in place, so input files written for 1.0 keep running on 1.1
  untouched. The list is pruned only when the interpretation of a key really
  has changed --- a renamed key, a changed default, a changed meaning. Dropping
  a series is therefore a deliberate act, it makes older input files fail
  loudly rather than silently produce a different answer, and every such
  removal is recorded in \cref{ch:revisions}.
\end{note}

The version alone does not identify a build. \file{JobInfo.txt} records the
version, the build date and the compiler, alongside the versions of the input
reader, grid and utility layers (\cref{sec:output:jobinfo}). A result directory
thus records which binary produced it \emph{within} one patch level as well as
across patch levels --- which is what makes a difference between two runs
traceable when both of them claim to be 1.0.0.

\section{What ``results may change'' means, and how it is bounded}
\label{sec:policy:results}

This is the section that matters. A solver whose numbers quietly drift between
releases is not usable for anything that has to be defended later, so the
commitment is this:

\begin{quote}
  \textbf{Within a major version, any change that alters numerical output is
  treated as a change of behavior. It is documented with a quantified
  difference --- the case it was measured on and how large the difference was
  --- and not merely mentioned.}
\end{quote}

That commitment is workable only because the repository already distinguishes
three situations, and keeps them apart deliberately.

\begin{description}

  \item[A refactoring, which must be bit-identical.] Restructuring that is not
    meant to change anything. The requirement is not ``small differences''; it is
    that \file{tools/compare.py} reads \code{0.000e+00} in every column of
    every compared case. This is achievable, rather than aspirational, because
    the solver is written to make it so --- the deterministic reductions, the
    CRS storage order, the registration order and the two direction orderings
    of \cref{ch:reproducibility} exist for exactly this reason, and the
    restructurings that produced the current code were all validated that way
    (\cref{sec:layout:workflow}).

  \item[A bug fix, which changes results because the old numbers were wrong.]
    Here the difference is the point, not an accident. The entry in
    \cref{ch:revisions} says what was wrong, which cases move, and by how much;
    if the old behavior was wrong only under a condition, the entry says which
    condition, so that you can tell whether your own results are affected. A
    bug fix of this kind is a PATCH release.

  \item[A deliberate change of method.] A different discretization, a different
    default, a different convergence test. These change results everywhere and
    are not bug fixes. They require at least a MINOR bump --- a MAJOR one if an
    existing input file would mean something different afterwards --- and an
    entry in \cref{ch:revisions} that quantifies the change on a named case.

\end{description}

The instrument is \file{tools/compare.py} in all three cases. It compares
\file{History.dat} and the final \file{field*.dat} column by column, normalizes
each column's largest difference by that column's own scale, and holds the four
solver-diagnostic columns --- \code{NK\_iter}, \code{NK\_resid},
\code{BCG\_iter}, \code{BCG\_resid} --- out of the pass/fail decision while
still reporting them. That last detail matters when you read a comparison: an
iteration count that moved while every physical column reads \code{0.000e+00}
is a change in how the answer was reached, not in the answer.

\begin{warn}[Compare like with like]
  A comparison across two linear-solver backends, or between a build with AMGCL
  and one without, cannot show bit-identity and is not evidence of anything:
  the backends agree only to the linear-solver tolerance
  (\cref{sec:linsolve:choosing}), and \val{"amgcl"} is not even reproducible
  under a change of thread count. Fix the backend and the thread count on both
  sides.
\end{warn}

\section{Input-file compatibility}
\label{sec:policy:compat}

\textbf{What is guaranteed.} An input file written for 1.0 runs on any 1.0.x,
and every key in it keeps the meaning this manual gives it. This covers a
block-structured \file{.inp} file as well as a TOML one, since the two produce the same
tree (\cref{sec:inputformat:legacy}). New optional keys may appear at 1.1; leaving them out
will continue to give you the behavior you have now.

\textbf{What is not guaranteed, and why it needs care.} A key that is renamed
or removed cannot go unnoticed: an entry the solver does not read stops the
run, naming it (\cref{sec:inputformat:unknown}), so an input file that uses an
old spelling fails loudly rather than running with that key at its default.
What no check can catch is a key that keeps its name and changes its default
or its meaning; the same goes for what a block left out of the file stands for
(\cref{sec:inputformat:defaults}).

\begin{warn}[A changed default is not diagnosed, so it is listed instead]
  Because the parser cannot warn you about it, the policy carries the burden
  instead: \textbf{a release that renames or removes a key, or changes a
  default, lists it in \cref{ch:revisions}}, with the old and the new
  spelling, or the old and the new value. Reading that list is the only
  defense against a changed default, so read it before running a case you
  cannot afford to re-verify on a new release.
\end{warn}

Keep the copy of the input file that every result directory contains
(\cref{sec:output:dir}). It is a verbatim copy of what produced that run, which
makes it the reference when you need to establish, later and under a newer
version, what a result was actually computed from.

\section{Output-format stability}
\label{sec:policy:outputs}

\file{field*.dat} is not only an output file: it is also the ASCII restart
format (\cref{sec:output:field}). Its column layout is therefore \textbf{fixed
within a major version} --- a restart file written by 1.0.0 is read by 1.0.x
and by 1.1, and the columns mean the same things. That is also why it is
written with 18 significant digits whatever \key{Output.Dimensioned} says.

\file{History.dat} and \file{CalcInfo.dat} are diagnostics, and the policy for
them is looser: \textbf{they may gain columns at the end in a minor release.}
Both files begin with a header line naming every column, so

\begin{quote}
  \textbf{read them by header name, never by column position.}
\end{quote}

A script that indexes \file{History.dat} by number will keep working until the
day a column is added and then quietly report the wrong quantity. A script that
looks the name up will not. Columns are not renamed and not reordered within a
major version; only appended.

\section{What is maintained and what is not}
\label{sec:policy:support}

Short and honest.

\begin{description}
  \item[Maintained.] The three models --- well-mixed in one dimension,
    well-mixed in two, and the through-thickness concentration model
    \val{"CZ"} (\cref{ch:concentration}) --- the shipped cases of
    \cref{ch:cases}, and the test suite of \cref{sec:verification:suite}. These
    are built, run and checked; a change that breaks one of them is a broken
    change.
  \item[Not promises.] Everything in \cref{ch:notverified}. Those items are
    known gaps in the evidence, listed so that you do not assume otherwise.
    They are not undertakings to close them, and an unverified quantity is not
    thereby a wrong one --- it is one for which Yuzen offers you nothing, and
    where you should supply your own evidence.
  \item[Not offered, and why.] An in-house BiCGStab backend with ILU(0),
    measured during development and left out of the release. In 1-D it was
    four to six times slower than \val{"eigen"}
    (\cref{sec:linsolve:choosing}), it needed 3676 linear iterations on
    \file{leveling\_1d\_reg} where \val{"eigen"} needs 21
    (\cref{sec:verification:backends}), and it had defects of its own; and
    \val{"eigen"} gives everything it offered, including bit-identity under
    a change of thread count
    (\cref{ch:reproducibility,sec:verification:linalg}). If you are choosing
    a backend for production work,
    \cref{sec:linsolve:choosing} is the section to read.
\end{description}

\section{Reporting a problem}
\label{sec:policy:reporting}

A report is actionable when somebody else can reproduce what you saw. Five
things make that possible, and they are all things you already have.

\begin{enumerate}
  \item \textbf{The version and the build date}, copied from
    \file{JobInfo.txt} in the result directory. ``1.0'' is not enough; the build
    date and the compiler distinguish two binaries at the same patch level.
  \item \textbf{The input file} --- the copy in the result directory, not the
    one in your working tree, which you may have edited since
    (\cref{sec:output:dir}).
  \item \textbf{The termination message}, the line near the end of the run
    log that says why the run ended (\cref{tab:termination}). The log goes to standard output and to no file,
    so it survives only if you redirected it (\cref{sec:output:running}). A
    run that stopped on \code{Steady state.} and a run that stopped on
    \code{Dry out.} are different reports.
  \item \textbf{Whether \file{Warnings.txt} and \file{FatalErrors.txt} have
    content}. Neither file carries a banner, so each is empty exactly when
    nothing was raised, and anything in either is a real message
    (\cref{sec:output:warnings}). A fatal error is written to
    \file{FatalErrors.txt} \emph{and} to standard error. The one abort that
    leaves no file is an input the solver rejects, which happens before the
    result directory exists; then the message is on standard error only.
  \item \textbf{The smallest case that still shows it.} Coarsen the mesh,
    shorten the end time, and switch off the physics terms one at a time until
    the behavior disappears; the term whose removal makes it go away is a
    large part of the report. A case that takes ten seconds gets looked at.
\end{enumerate}

\begin{note}[A thread-count difference is always a bug]
  If a case gives different results at different thread counts on the
  \val{"eigen"} backend, that is a defect by definition
  --- thread invariance is one of the claims the test suite asserts
  (\cref{ch:reproducibility,sec:verification:linalg}) --- and it is worth
  reporting as such, with the two \file{compare.py} outputs attached. The one
  case where it is expected rather than wrong is \val{"amgcl"}, whose internal
  reductions are non-deterministic (\cref{sec:linsolve:choosing}).
\end{note}
